% Generated draft: requires mathematical and source review.
\documentclass[11pt]{article}
\usepackage[margin=25mm]{geometry}
\usepackage{fontspec}
\setmainfont{Latin Modern Roman}
\usepackage{amsmath,amssymb,mathtools,unicode-math}
\setmathfont{Latin Modern Math}
\usepackage{enumitem,needspace,xurl,hyperref}
\setlength{\parindent}{0pt}
\setlength{\parskip}{5pt}
\setlength{\emergencystretch}{3em}
\newcommand{\N}{\mathbb{N}}
\newcommand{\Z}{\mathbb{Z}}
\newcommand{\Q}{\mathbb{Q}}
\newcommand{\R}{\mathbb{R}}
\newcommand{\C}{\mathbb{C}}
\newcommand{\F}{\mathbb{F}}
\newcommand{\A}{\mathbb{A}}
\newcommand{\PP}{\mathbb{P}}
\newcommand{\E}{\mathbb{E}}
\newcommand{\eps}{\varepsilon}
\newcommand{\dd}{\,\mathrm{d}}
\begin{document}
\section*{UnsolvedMath ID: 30006052}
\textbf{Source problem code:} catalogue entry
\subsection{Definitions and mathematical statement}
Let \(\mathscr K\) denote the set of knot types in \(S^3\) under conventions to be specified by the random-knot model. A knot property \(\mathcal P\) is represented by a subset \(A_{\mathcal P}\subseteq\mathscr K\). For each index \(n\), let \((\Omega_n,\mathcal F_n,\Pr_n)\) be a probability space and let \(K_n\) be a random knot taking values in \(\mathscr K\), with \(\{\omega\in\Omega_n:K_n(\omega)\in A_{\mathcal P}\}\in\mathcal F_n\). Define
\[
p_n(\mathcal P)=\Pr_n\!\left(K_n\in A_{\mathcal P}\right).
\]
The spaces need not be coupled across different indices. The knot conventions, distributions, meaning of the size parameter \(n\), and admissible indices are part of the model; the supplied record does not fully specify them.

\textbf{Parametric problem.} For a fixed property \(\mathcal P\) and a specified random-knot model \((K_n)\), determine
\[
L_{\mathcal P,(K_n)}=\lim_{n\to\infty}p_n(\mathcal P)
\]
when this limit exists. Existence is a qualification of the question, not an assertion for every property and every model. The original excerpt asks for the probability that a random knot satisfies a property arising from preceding open problems; those preceding properties are not identified in the supplied excerpt. The asymptotic formulation above follows the supplied cleaned statement and the excerpt's discussion of increasing grid size. Neither a particular property nor a complete probability model is fixed, so this is a research question parameterized by those choices, rather than a single specified numerical target.

\textbf{Grid-model context and unresolved qualifications.} The excerpt describes choosing a random grid diagram of size \(n\) and studying the represented knot as \(n\to\infty\). It states that every knot in \(S^3\) has a grid representation and gives the finite count \(G(n)=n!(n-1)!\), while describing every counted diagram as representing a knot. These are supplied source claims: the precise diagram-counting conventions, treatment of diagrams representing links with multiple components, admissible small sizes, and sampling distribution are not supplied. Consequently, this draft does not adopt that count as a definition of a fully specified probability model, assume uniform sampling, or identify probabilities on diagrams with uniform probabilities on knot types. The excerpt also reports computer experiments on invariants related to smooth sliceness using GridPyM; it supplies neither a precise sliceness probability question nor a limiting conclusion. Its references to contact topology, knot Floer homology, and monotonic simplification of unknot grid diagrams provide background, not additional hypotheses.
\subsection{Short English statement}
For a specified way of sampling knots of increasing size and a fixed knot property, determine the limiting probability of that property, if the limit exists. The supplied record leaves the property and probability model open. Its grid-diagram example requires further counting and sampling conventions before it defines a precise instance of the question.
\subsection{Sources}
\begin{itemize}
\item Local UnsolvedMath registry, supplied complete record for permanent ID \(30006052\), title \emph{Probabilities of Properties in Random-Knot Models}, source problem code \texttt{catalogue entry}, in the Oberwolfach Reports Open Problems collection. The registry supplies the labels \texttt{open} and \texttt{corrected\_verified}; these are supplied curation labels, not independent mathematical or source verification by this draft. Its verification note states that the official report was checked and the question isolated and its notation repaired; that statement is preserved as registry provenance only.
\item Supplied primary-source citation: \emph{MATRIX-MFO Tandem Workshop: Invariants in Low-Dimensional Topology: Combinatorics, Geometry, and Computation} (2025), \emph{Oberwolfach Reports}, pp.\ 2511--2550, DOI \url{https://doi.org/10.4171/owr/2024/43}. The supplied literature list calls this \emph{Oberwolfach Report 43 (2024)}. The differing year designations are retained without resolution.
\item The supplied excerpt places the grid-model discussion under \emph{Models for random knots}, followed by the names Valentina Bais, Agnese Barbensi, Benjamin Burton, Daniele Celoria, Chun-Sheng Hsueh, Marc Kegel, Joan Licata, Alice Merz, Jonathan Spreer, Annika Thiele, and Lucy Tobin. The registry does not identify an individual proposer; this draft does not assign that role to any listed person.
\item Supplied literature source: \emph{Random Meander Model for Links}, \url{https://doi.org/10.1007/s00454-024-00663-9}. The record supplies no theorem from this reference that resolves the general parametric question.
\item Literature assessment dated exactly \texttt{2026-08-22}: the supplied curation labels the question open, describes it as model-dependent and a direction for future probabilistic study, and reports that no general probability theorem for arbitrary knot properties was found. It lists the two DOI sources above as checked. This is supplied dated curation, not an independently verified literature assessment; the record requests independent verification of time-sensitive status and cited claims.
\item The excerpt mentions GridPyM as reference [3], monotonic simplification of unknot grid diagrams as [13], contact topology as [22], and knot Floer homology as [23]. Their bibliographic entries are absent from the supplied record, so those references cannot be identified reliably here.
\end{itemize}
\end{document}
